\section{Approach}

\subsection{Time-Expanded Network and DP Formulation}

The assembly campaign is modeled as a finite-horizon Markov
decision process over a time-expanded network. Decision epochs are
discretized into configurable periods of length $\Delta t$
(default: 7 days). The system state at epoch $t$ is

\begin{equation}
  S(t) = \bigl\{\mathcal{M}(t),\; \mathcal{C}(t),\;
                \mathcal{Q}(t),\; n_{\mathrm{tugs}}(t)\bigr\}
\end{equation}

\noindent where $\mathcal{M}(t)$ is the set of completed modules,
$\mathcal{C}(t)$ is the set of docked crew vehicles with their
remaining rotation periods, $\mathcal{Q}(t)$ is the cargo-in-transit
queue with arrival times, and $n_{\mathrm{tugs}}(t)$ is the number
of reusable transfer stages available at LEO.

At each epoch the optimizer selects a decision
$D(t) = \{\text{launches},\;\text{crew allocation},\;
\text{assembly tasks},\;\text{tug assignments}\}$
that transitions the system to state $S'$. The Bellman optimality
equation is

\begin{equation}
  V(S,t) = \min_{D \in \mathcal{D}(S,t)}
            \bigl[\, c(S, D, t) + V(S', t+1) \,\bigr]
\end{equation}

\subsection{Multi-Objective Cost Function}

Three objectives are minimized simultaneously via a weighted sum
with safe normalization:

\begin{equation}
  J = w_1 \frac{N_{\mathrm{launches}}}{N_{\mathrm{launches}}^{\max}}
    + w_2 \frac{T}{T^{\max}}
    + w_3 \frac{C}{C^{\max}}
  \label{eq:cost}
\end{equation}

\noindent where weights $w_1, w_2, w_3 \geq 0$ are user-adjustable.
Upper bounds are $N_{\mathrm{launches}}^{\max} = 2|\mathcal{M}_{\mathrm{total}}|$,
$T^{\max} = T_{\mathrm{horizon}}$, and
$C^{\max} = \$2{,}000|\mathcal{M}_{\mathrm{total}}|$\,M.
Safe normalization returns 0 when the upper bound is zero.

\subsection{Transfer Model}

Payload mass delivered to L4 is computed via the Tsiolkovsky
rocket equation applied to the LEO-to-L4 transfer:

\begin{equation}
  m_{\mathrm{delivered}} = m_{\mathrm{LEO}}
    \cdot \exp\!\left(\frac{-\Delta v}{I_{sp}\, g_0}\right)
  \label{eq:rocket}
\end{equation}

\noindent with $g_0 = 9.807 \times 10^{-3}$\,km/s$^2$. Two
transfer options are modeled: a direct Hohmann-like transfer
($\Delta v = 4.1$\,km/s, 60-day transit) and a low-energy
weak-stability-boundary (WSB) transfer ($\Delta v = 3.8$\,km/s,
120-day transit)~\cite{belbruno_wbs}.

For LEO-only vehicles that require a transfer stage of dry mass
$m_d$ and propellant mass $m_p$, the tug mass ratio $R =
\exp(\Delta v / I_{sp} g_0)$ constrains maximum deliverable cargo:

\begin{align}
  m_{\mathrm{cargo,max}} &= \frac{R\,m_d - (m_d + m_p)}{1 - R}
  \label{eq:tug} \\
  m_{\mathrm{delivered}} &= \min\!\left(m_{\mathrm{payload}},\;
                               m_{\mathrm{cargo,max}}\right)
\end{align}

\subsection{Crew Model}

All crew across docked vehicles form a single shared pool:

\begin{equation}
  n_{\mathrm{crew}} = \sum_{v \in \mathcal{C}(t)} n_v
\end{equation}

EVA operations require a buddy system (pairs) and intravehicular
activity (IVA) monitors. The stable allocation is found
iteratively:

\begin{equation}
  n_{\mathrm{IVA}} = \left\lceil
    \frac{n_{\mathrm{EVA\,pairs}}}{k_{\mathrm{IVA}}}
  \right\rceil, \quad
  n_{\mathrm{EVA\,pairs}} = \left\lfloor
    \frac{n_{\mathrm{crew}} - n_{\mathrm{IVA}}}{2}
  \right\rfloor
  \label{eq:eva}
\end{equation}

\noindent where $k_{\mathrm{IVA}}$ is the maximum EVA pairs per
IVA monitor (default: 2). Productive assembly hours per period
are reduced by the duty cycle $\eta_d = 0.6$ to account for
rest and maintenance overhead:

\begin{equation}
  H_{\mathrm{EVA}} = n_{\mathrm{EVA\,pairs}} \cdot h_{\mathrm{session}}
                     \cdot d_{\mathrm{period}} \cdot \eta_d
  \label{eq:crew_hours}
\end{equation}

Robotic arm stations operate continuously (24/7) but with a
time penalty $\eta_r = 1.5$ relative to crew EVA:

\begin{equation}
  H_{\mathrm{robotic}} = \frac{n_{\mathrm{arms}} \cdot h_{\mathrm{period}}}{\eta_r}
\end{equation}

Each crew vehicle trades unused crew capacity for cargo mass:

\begin{equation}
  m_{\mathrm{cargo}} = m_{\mathrm{capacity}}
                       - n_{\mathrm{crew}} \cdot m_{\mathrm{per\,crew}}
\end{equation}

\subsection{Proximity Operations and Congestion Model}

Simultaneous vehicle presence at the assembly site increases
coordination overhead, reducing effective assembly throughput.
A scale-adaptive penalty function is applied to available crew
and robotic hours each period:

\begin{equation}
  \pi(n,\,p) = 1 + \frac{\alpha\,(n-1)^{\beta}}{C(p)}, \quad
  C(p) = C_{\mathrm{base}} + (C_{\mathrm{max}} - C_{\mathrm{base}})\,p
  \label{eq:penalty}
\end{equation}

\noindent where $n$ is the number of vehicles in proximity
(crew vehicles plus arriving cargo vehicles), $p =
|\mathcal{M}_{\mathrm{built}}| / |\mathcal{M}_{\mathrm{total}}|$
is fractional build progress, and default parameters are
$\alpha = 0.1$, $\beta = 1.5$, $C_{\mathrm{base}} = 2$,
$C_{\mathrm{max}} = 10$. Effective assembly hours are:

\begin{equation}
  H_{\mathrm{EVA,eff}} = \frac{H_{\mathrm{EVA}}}{\pi(n,p)}, \quad
  H_{\mathrm{robotic,eff}} = \frac{H_{\mathrm{robotic}}}{\pi(n,p)}
\end{equation}

Cumulative collision risk is tracked as a safety metric:

\begin{equation}
  \mathcal{R} = \sum_{t=0}^{T} \kappa \cdot \binom{n(t)}{2} \cdot \Delta t
  \label{eq:risk}
\end{equation}

\subsection{Partial Assembly (Work-in-Progress)}

Modules whose assembly hours exceed one period's capacity are
carried forward as work-in-progress (WIP). Remaining hours
$h_r$ are updated each period:

\begin{equation}
  h_r^{(t+1)} = \max\!\left(0,\;
    h_r^{(t)} - \min\!\left(h_r^{(t)},\; H_{\mathrm{available}}^{(t)}\right)
  \right)
\end{equation}

\noindent A module is marked complete when $h_r \leq \varepsilon$
(default $\varepsilon = 0.01$\,h).

\subsection{Pad Turnaround Constraint}

Each launch vehicle $v$ is subject to a minimum inter-launch
interval to model launch pad refurbishment:

\begin{equation}
  t - t_{\mathrm{last}}(v) \geq \Delta t_{\mathrm{pad}},
  \quad \forall\, v \in \mathcal{V}_{\mathrm{cargo}} \cup
  \mathcal{V}_{\mathrm{crew}}
\end{equation}

\noindent Default: $\Delta t_{\mathrm{pad}} = 2$ periods (14 days).

\subsection{Beam Search Pruning}

Exact DP over the assembled module bitmask is intractable for
structures with $|\mathcal{M}_{\mathrm{total}}| \geq 20$
($2^{50}$ states for a 50-module spacecraft). Forward beam search
retains only the top-$K$ states per epoch ranked by the
progress heuristic:

\begin{equation}
  \phi(S) = 10\,|\mathcal{M}_{\mathrm{built}}|
          + 3\,|\mathcal{D} \setminus \mathcal{M}_{\mathrm{built}}|
          + |\mathcal{Q}|
          + (5 + 3r)\cdot\mathbf{1}[n_{\mathrm{crew}}>0]
          + 0.5\,|\mathcal{W}|
  \label{eq:phi}
\end{equation}

\noindent where $r = |\mathrm{avail}(\mathcal{M}_{\mathrm{built}})
\cap \mathcal{D}|$ is modules ready to assemble on-site, and
$\mathcal{W}$ is the current WIP set. Default beam width:
$K = 100$.
